\section{Tensor geometry from singular directions}
\label{sec:geometry}

We now construct the PPT pair in Theorem~\ref{thm:main-pair}.  The geometric
step uses a map on $\Mat4$ with twenty specified singular outputs.  It turns
each singular output into a product vector in the kernel of a Choi matrix
on $\C^{10}\otimes\C^{10}$.  A condition on quadratic evaluations then
excludes product vectors from the range.

For a vector $x$, write $P_x=xx^*$, without normalizing it.  A projective
direction $[x]$ is the complex line spanned by a nonzero vector $x$.

\begin{proposition}[Tensor criterion]\label{prop:tensor-criterion}
Let $L:\Mat4\to\Mat4$ be PPT.  Suppose that
$\mathcal X=\{x_1,\ldots,x_{20}\}\subset\C^4\setminus\{0\}$ represents
twenty distinct projective directions and has the following properties:
\begin{enumerate}
\item $\rank L(P_x)<4$ for every $x\in\mathcal X$, and every nonzero
homogeneous quadratic in $\C[x_0,x_1,x_2,x_3]$ vanishes on at most nine
members of $\mathcal X$;
\item $L(P_y)$ is positive definite for some $y\in\C^4$.
\end{enumerate}
Then there are PPT maps $\Phi_1,\Phi_2:\Mat{10}\to\Mat{10}$, constructed
explicitly from $L$, such that
\[
 Z=J(\Phi_2\circ\Phi_1)\ne0,\qquad
 \ran Z\cap\{u\otimes v:u,v\in\C^{10}\}=\{0\}.
\]
In particular, $\Phi_2\circ\Phi_1$ is not entanglement breaking.
\end{proposition}

Section~\ref{sec:certificate} supplies an integer construction of $L$ with
these properties.  We prove the criterion by using the symmetric square to
record quadratic monomials and the exterior square to detect rank loss.

\begin{proof}
\emph{Symmetric and exterior tensors.}
Index an orthonormal basis of $\C^{10}$ by pairs $0\leq i\leq j<4$, and an
orthonormal basis of $\C^6$ by pairs $0\leq i<j<4$, both in lexicographic
order.  Define isometries $U:\C^{10}\to\C^4\otimes\C^4$ and
$V:\C^6\to\C^4\otimes\C^4$ by
\begin{align}
 Ue_{ii}&=e_i\otimes e_i,&
 Ue_{ij}&=\frac{e_i\otimes e_j+e_j\otimes e_i}{\sqrt2}\quad(i<j),
 \notag\\
 Ve_{ij}&=\frac{e_i\otimes e_j-e_j\otimes e_i}{\sqrt2}\quad(i<j).
 \label{eq:UV}
\end{align}
Their ranges are the symmetric and antisymmetric tensor subspaces.  Put
\begin{equation}
 \widehat x=U^*(x\otimes x)
 =\sum_i x_i^2e_{ii}+\sqrt2\sum_{i<j}x_ix_je_{ij}
 \qquad(x\in\C^4).
 \label{eq:veronese}
\end{equation}
Thus $U\widehat x=x\otimes x$ and
$UP_{\widehat x}U^*=P_x\otimes P_x$.  The coordinates of $\widehat x$ are
the ten quadratic monomials with nonzero normalization factors.  Hence the
quadratic hypothesis is equivalent to the assertion that any ten vectors
among $\{\widehat x:x\in\mathcal X\}$ span $\C^{10}$: a failure to span
would give a nonzero vector $u$ with $u^*\widehat x=0$ at those ten points,
and $u^*\widehat x$ is a nonzero homogeneous quadratic.

Identify $\C^6$ with $\bigwedge^2\C^4$ by $e_{ij}=e_i\wedge e_j$, with this
basis orthonormal.  Under this identification,
\[
 V(a\wedge b)=\frac{a\otimes b-b\otimes a}{\sqrt2}.
\]
For $i<j$, let $k<l$ be the complementary indices, and let
$\epsilon_{ijkl}$ be the sign of the ordering $(i,j,k,l)$ of $(0,1,2,3)$.
Define the real signed permutation $K$ on $\C^6$ by
\begin{equation}
 Ke_{ij}=\epsilon_{ijkl}e_{kl}.
 \label{eq:complement}
\end{equation}
Explicitly,
\[
 (Ke_{01},Ke_{02},Ke_{03},Ke_{12},Ke_{13},Ke_{23})
 =(e_{23},-e_{13},e_{12},e_{03},-e_{02},e_{01}),
\]
so $K=K^*=K^{-1}$.  The conjugate operation $w\mapsto K\overline w$ is
the complementary-pair exterior dual also used by Schliemann, Loss, and
MacDonald for two fermions in a four-dimensional one-particle space
\cite[Section III, Equations (26)--(28)]{SchliemannLossMacDonald2001}.

Define $S,R:\Mat{10}\to\Mat6$ by
\begin{equation}
 S(A)=V^*(L\otimes L)(UAU^*)V,\qquad
 R(A)=K\,S(A)^{\T}K^*.
 \label{eq:SR}
\end{equation}
The embeddings and compressions are CP, so Lemma~\ref{lem:ppt-stability}
makes $S$ PPT.  The map $T_6\circ S$ is CP by this PPT property, and its
output transpose is $S$, which is CP.  Thus $T_6\circ S$ is PPT, and the
CP conjugation by $K$ makes $R$ PPT.

\emph{The composite Choi matrix.}
The rectangular maps give an endomorphism
$R^\dagger\circ S\circ T_{10}$ of $\Mat{10}$.  To express it as the
composition of two maps on that algebra, let
$E:\C^6\to\C^{10}$ be the real coordinate inclusion into the first six
positions and define
\begin{equation}
 \Phi_1(A)=E\,S(A^{\T})E^*,\qquad
 \Phi_2(B)=R^\dagger(E^*BE).
 \label{eq:phi-pair}
\end{equation}
For $\Phi_1$, Lemma~\ref{lem:ppt-stability} makes $S\circ T_{10}$ CP by
input/output transpose equivalence, and makes its output transpose
$T_6\circ S\circ T_{10}$ CP by simultaneous transposition of $S$.
Thus $S\circ T_{10}$ is PPT, and the CP embedding by $E$ preserves that
property.  The adjoint and CP-composition clauses of the same lemma make
$\Phi_2$ PPT.  The formulas determine both maps explicitly once $L$ is
given.  For example, if
$L(A)=\sum_r Q_rAQ_r^*$, the adjoints in the construction are
\begin{align*}
 L^\dagger(Y)&=\sum_r Q_r^*YQ_r,\\
 S^\dagger(Y)&=U^*(L^\dagger\otimes L^\dagger)(VYV^*)U,\\
 R^\dagger(Y)&=S^\dagger\bigl((K^*YK)^{\T}\bigr).
\end{align*}

Since $E^*E=I_6$,
$\Phi_2\circ\Phi_1=R^\dagger\circ S\circ T_{10}$.  Its Choi matrix $Z$ is
positive.  The convention \eqref{eq:choi-definition} gives, for arbitrary
matrices $a,b$ of the appropriate sizes,
\[
 \tr\bigl[(a\otimes b)J(F)\bigr]=\tr\bigl[bF(a^{\T})\bigr].
\]
For Hermitian $a,b\in\Mat{10}$, the Hilbert--Schmidt adjoint identity
therefore gives
\begin{equation}
 \tr\bigl[(a\otimes b)Z\bigr]
 =\tr\bigl[bR^\dagger(S(a))\bigr]
 =\tr\bigl[R(b)S(a)\bigr].
 \label{eq:choi-pairing}
\end{equation}
The Choi transpose has canceled the input transpose in
\eqref{eq:phi-pair}.  We will prove that the right side vanishes at
$a=b=P_{\widehat x}$ for each singular direction.  Positivity of $Z$ will
then give the kernel vector $\widehat x\otimes\widehat x$.
The positive definite output will separately show that $Z$ is nonzero.

\emph{Singular outputs and kernel products.}
For any operator $B$ on $\C^4$,
\[
 V^*(B\otimes B)V=\bigwedge\nolimits^2B.
\]
Both sides send $a\wedge b$ to $Ba\wedge Bb$.  If
$B_x=L(P_x)\succeq0$ and $W_x=\ran B_x$, it follows that
\begin{equation}
 S(P_{\widehat x})=\bigwedge\nolimits^2B_x,\qquad
 \ran S(P_{\widehat x})=\bigwedge\nolimits^2W_x.
 \label{eq:exterior-range}
\end{equation}
Indeed, in an orthonormal eigenbasis of $B_x$, the exterior eigenvalues are
$\lambda_i\lambda_j$ for $i<j$.  This proves both positivity and the range
identity.  A Hermitian matrix $D$ satisfies $D^{\T}=\overline D$, so
transposition conjugates its range.  Consequently
\begin{equation}
 \ran R(P_{\widehat x})=K\overline{\bigwedge\nolimits^2W_x}.
 \label{eq:complement-range}
\end{equation}

For exterior two-vectors $w,u$, let $[w\wedge u]_{0123}$ denote the
coefficient of $e_0\wedge e_1\wedge e_2\wedge e_3$.  With the Hermitian
inner product conjugate-linear in its first argument, the definition of $K$
gives
\begin{equation}
 \langle K\overline w,u\rangle
 =\sum_{i<j}w_{ij}\epsilon_{ijkl}u_{kl}
 =[w\wedge u]_{0123}.
 \label{eq:wedge-pairing}
\end{equation}
Here $k<l$ is complementary to $i<j$ in each term.  If
$w,u\in\bigwedge^2W$ and $\dim W<4$, then
$w\wedge u\in\bigwedge^4W=0$.  Equations
\eqref{eq:exterior-range}--\eqref{eq:wedge-pairing} show that
$S(P_{\widehat x})$ and $R(P_{\widehat x})$ have orthogonal ranges for every
$x\in\mathcal X$.  This includes $\dim W_x<2$, when both exterior ranges
are zero.  The conjugation in \eqref{eq:wedge-pairing} is essential for this
conclusion over complex subspaces.

For $x\in\mathcal X$, orthogonality of the two positive output ranges gives
\[
 (\widehat x\otimes\widehat x)^*Z(\widehat x\otimes\widehat x)
 =\tr\bigl[R(P_{\widehat x})S(P_{\widehat x})\bigr]=0.
\]
For a positive matrix, $v^*Zv=0$ implies $Zv=0$.  Hence
\begin{equation}
 \widehat x\otimes\widehat x\in\ker Z\qquad(x\in\mathcal X).
 \label{eq:kernel-vectors}
\end{equation}
For the input $y$ in the second hypothesis,
$S(P_{\widehat y})=\bigwedge^2(L(P_y))$ is positive definite, and so is
$R(P_{\widehat y})=K S(P_{\widehat y})^{\T}K^*$.  The trace of their
product is strictly positive.  Equation~\eqref{eq:choi-pairing} at
$a=b=P_{\widehat y}$ therefore proves $Z\ne0$.

\emph{Excluding product vectors from the range.}
It remains to exclude a product vector from the range.
The following zero-set argument uses the product-exclusion principle of
Bennett and coauthors \cite[Lemma~1]{Bennett1999}, in the nonorthogonal
form given by Pittenger \cite[Section~IV, Lemma~3]{Pittenger2003}, and
is related to Parthasarathy's polynomial interpolation construction
\cite[Proposition~1.3]{Parthasarathy2004}.
Suppose that a nonzero product $u\otimes v$ belongs to $\ran Z$.  Since
$Z$ is Hermitian, its range is the orthogonal complement of its kernel.
Equation~\eqref{eq:kernel-vectors} implies
\[
 p_u(x)p_v(x)=0\quad(x\in\mathcal X),\qquad
 p_u(x)=u^*\widehat x,\quad p_v(x)=v^*\widehat x.
\]
For example,
\[
 p_u(x)=\sum_i\overline{u_{ii}}x_i^2
       +\sqrt2\sum_{i<j}\overline{u_{ij}}x_ix_j.
\]
This is a homogeneous quadratic in $x$; only its fixed coefficients are
conjugated.  Its distinct monomials make it nonzero when $u\ne0$.
The same holds for $p_v$, since $v\ne0$.  Each polynomial therefore
vanishes on at most nine members of $\mathcal X$.  Their zero sets cannot
cover all twenty members, because $9+9<20$.  This contradiction proves
that $\ran Z$ contains no nonzero product vector.

By Lemma~\ref{lem:product-range}, the nonzero positive matrix $Z$ is
nonseparable.  Lemma~\ref{lem:choi-separable} then shows that the composite
is not entanglement breaking.
\end{proof}
